% Figure and table blocks for the Sagnac report (prepared 22 October 2026, before the first draft;
% numbers of 23 October, after the calibration file was completed).
%
% Numbers only. Every value is from sagnac.txt and calibration.txt (sagnac.py, calibration.py), the
% final analysis of RESULTS.md, rounded by the course convention (the uncertainty to one significant
% figure, the value to the same digit). There is no report prose in this file: each \caption{} holds
% only a visible placeholder and, in a comment, what the caption has to say (the same list as in
% REPORT_TODO.md, "Figures and tables"). Write the captions and all text in your own words.
%
% Needs graphicx, siunitx, booktabs and cleveref, as in the Chua report, and the figures on the
% graphics path:   \graphicspath{{report_figures/}}
% A figure* (full text width) never appears on the page where it is defined: put its block a
% paragraph or two before the first text that cites it (the Chua report's Fig. 7 lesson).
%
% Unit of the Sagnac slope, defined once. rad/(rad/s) is a second: change this to \unit{\second}
% if you prefer, and the tables follow.
\providecommand{\slopeunit}{\ensuremath{\mathrm{rad}/(\mathrm{rad\,s^{-1}})}}


% ===== Figures =============================================================================

% ---- Fig. 1, section 1.2 (setup) ---------------------------------------------------------
\begin{figure}[tb]
  \centering
  \includegraphics[width=\columnwidth]{setup.pdf}
  \caption{[caption]%
    % Must say: top view of the apparatus, after Fig. 1 of the manual (cite it); P1, P2 polarizers,
    % L1, L2 coupling lenses into the two ends of the fiber, L3 the camera lens; beam 1 and beam 2
    % travel the same loop (N turns of radius R) in opposite directions and meet again at the beam
    % splitter; everything drawn turns with the table, about the loop's center; the sense of
    % rotation was not recorded, so the omega arrow is schematic; the small off-center loop of
    % task 4 is not drawn.
  }
  \label{fig:setup}
\end{figure}

% ---- Fig. 2, section 2.1 (rotation calibration) -------------------------------------------
\begin{figure}[tb]
  \centering
  \includegraphics[width=\columnwidth]{calibration.pdf}
  \caption{[caption]%
    % Must say: omega of the table at each controller setting S, from the timed turns (error bars:
    % the timing uncertainty, mostly smaller than the markers); line: the linear calibration used
    % for the photographs; band: its +-1 sigma uncertainty, which is the uncertainty of omega at any
    % photo setting. Fit results in the text, not here.
  }
  \label{fig:calibration}
\end{figure}

% ---- Fig. 2b, section 2.1 or 2.3 (optional: the photo settings converted to omega) ------------
% The same numbers as the omega column of Table 2: use the figure or the table, not both.
\begin{figure}[tb]
  \centering
  \includegraphics[width=\columnwidth]{conversion.pdf}
  \caption{[caption]%
    % Must say: the nine photo settings read off the line of Fig. 2: dotted guides from each
    % setting to the line and across to the omega axis, the bar at the line its uncertainty.
    % Nothing was measured at these settings; the values are those of Table 2.
  }
  \label{fig:conversion}
\end{figure}

% ---- Fig. 3, section 2.2 (the patterns, full width) ---------------------------------------
\begin{figure*}[tb]
  \centering
  \includegraphics[width=\textwidth]{patterns.pdf}
  \caption{[caption]%
    % Must say: one photograph per setting, at rest and at the nine speeds in the order taken; the
    % fitted dark rings (dashed) and ring center (+); omega from the calibration. The rings move
    % inward as omega rises; the lower left has no fringes.
  }
  \label{fig:patterns}
\end{figure*}

% ---- Fig. 4, section 2.2 (rings evenly spaced in r^2) -------------------------------------
\begin{figure}[tb]
  \centering
  \includegraphics[width=\columnwidth]{rings.pdf}
  \caption{[caption]%
    % Must say: intensity averaged around the ring center for the rest photo and the top-speed
    % photo of Fig. 3, (a) against r, (b) against r^2; dashed: the dark-ring radii; arrows: the
    % top-speed spacings, unequal in r, equal in r^2 (P, Eq. rings).
  }
  \label{fig:rings}
\end{figure}

% ---- Fig. 5, section 2.3 (phase of one photo) ---------------------------------------------
\begin{figure}[tb]
  \centering
  \includegraphics[width=\columnwidth]{extraction.pdf}
  \caption{[caption]%
    % Must say: (a) the top-speed photo with the located dark-ring minima (one color per ring)
    % and the circles fitted to them (dashed), center +; (b) R_m^2/P against the ring number m at
    % rest and at the top speed with the fitted lines (Eq. dark rings); the intercepts differ by
    % 0.52 fringe, i.e. Delta phi = 2 pi x 0.52 = 3.3 rad.
  }
  \label{fig:extraction}
\end{figure}

% ---- Fig. 6, section 2.3 (the result) -----------------------------------------------------
\begin{figure}[tb]
  \centering
  \includegraphics[width=\columnwidth]{phase_vs_omega.pdf}
  \caption{[caption]%
    % Must say: the phase shift of each setting (ring minima) against omega from the calibration;
    % error bars: the calibration in omega, the standard error of the photos in Delta phi; solid:
    % the fitted line; dashed with band: the Sagnac prediction a_th omega (Eq. sagnac) with its
    % 1 sigma band from N and R.
  }
  \label{fig:phase}
\end{figure}

% ---- Fig. 7, section 2.4 (phase-reading methods) ------------------------------------------
\begin{figure}[tb]
  \centering
  \includegraphics[width=\columnwidth]{methods.pdf}
  \caption{[caption]%
    % Must say: (a) the phase shift by the cosine fit and by template matching, minus the
    % ring-minima value; (b) the slope from each method with its fit uncertainty; band: the
    % prediction +-1 sigma.
  }
  \label{fig:methods}
\end{figure}

% ---- Fig. 8, section 2.6 (drift) ----------------------------------------------------------
\begin{figure}[tb]
  \centering
  \includegraphics[width=\columnwidth]{drift.pdf}
  \caption{[caption]%
    % Must say: (a) omega against the time since the first photo: the settings were taken in
    % increasing order; (b) each photo's phase shift minus the line of Fig. 6 against time, with a
    % trend line per setting: the phase wanders within and between settings.
  }
  \label{fig:drift}
\end{figure}

% ---- optional, section 2.4 (random pixels): include only if the paragraph on prep Q5 uses it -
% \begin{figure}[tb]
%   \centering
%   \includegraphics[width=\columnwidth]{random_pixels.pdf}
%   \caption{[caption]%
%     % Must say: (a) the slope from ten fixed pixels over 1000 random draws (outline: the same
%     % pixels with each photo's own center and P; lines: ring-minima slope and prediction);
%     % (b) the spread between draws, and the fit error of one draw, against the number of pixels
%     % per draw.
%   }
%   \label{fig:pixels}
% \end{figure}


% ===== Tables ==============================================================================
% \small and a narrower column gap keep every table inside its column; drop them if your template
% already sets table text smaller.

% ---- Table 1, section 2.1 (optional: the timed settings; Fig. 2 and the text can carry them) --
\begin{table}[tb]
  \centering\small\setlength{\tabcolsep}{4pt}
  \begin{tabular}{cccc}
    \toprule
    $S$ & $n$ & $\langle T\rangle$ (\unit{\second}) & $\omega$ (\unit{\radian\per\second}) \\
    \midrule
    5.78  & 5  & \num{65.45\pm0.09} & \num{0.0960\pm0.0001} \\
    14.67 & 5  & \num{22.33\pm0.07} & \num{0.2813\pm0.0009} \\
    25.66 & 5  & \num{12.42\pm0.06} & \num{0.506\pm0.002} \\
    36.63 & 5  & \num{8.60\pm0.06}  & \num{0.730\pm0.005} \\
    45.09 & 5  & \num{6.92\pm0.06}  & \num{0.908\pm0.008} \\
    56.03 & 10 & \num{5.57\pm0.03}  & \num{1.128\pm0.006} \\
    68.85 & 10 & \num{4.56\pm0.03}  & \num{1.38\pm0.01} \\
    79.21 & 10 & \num{3.99\pm0.03}  & \num{1.58\pm0.01} \\
    90.00 & 10 & \num{3.55\pm0.03}  & \num{1.77\pm0.02} \\
    \bottomrule
  \end{tabular}
  \caption{[caption]%
    % Must say: n consecutive turns timed at each controller setting S; <T> the mean turn time
    % with its uncertainty (spread of the turns and the 0.2 s reaction time of the presses);
    % omega = 2 pi/<T>.
  }
  \label{tab:calibration}
\end{table}

% ---- Table 2, section 2.3 (the measured values per setting) --------------------------------
\begin{table}[tb]
  \centering\small\setlength{\tabcolsep}{3.5pt}
  \begin{tabular}{lccccc}
    \toprule
    $S$ & Photos & First photo & $\omega$ & $\Delta\phi$ & $P$ \\
        &        &             & (\unit{\radian\per\second}) & (\unit{\radian}) & ($10^5$\,px$^2$) \\
    \midrule
    rest & 10 & 16:21:40 & 0 & 0 & \num{1.075}\\
    12.59 & 15 & 16:24:10 & \num{0.236\pm0.003} & \num{0.39\pm0.03} & \num{1.071}\\
    20.74 & 15 & 16:25:44 & \num{0.403\pm0.005} & \num{0.64\pm0.02} & \num{1.058}\\
    31.11 & 19 & 16:26:38 & \num{0.616\pm0.008} & \num{1.02\pm0.02} & \num{1.065}\\
    44.07 & 22 & 16:27:54 & \num{0.88\pm0.01} & \num{1.59\pm0.02} & \num{1.093}\\
    55.83 & 18 & 16:29:08 & \num{1.12\pm0.01} & \num{1.984\pm0.008} & \num{1.058}\\
    64.00 & 16 & 16:30:16 & \num{1.29\pm0.02} & \num{2.375\pm0.008} & \num{1.094}\\
    73.65 & 16 & 16:31:28 & \num{1.49\pm0.02} & \num{2.700\pm0.004} & \num{1.087}\\
    82.18 & 15 & 16:32:32 & \num{1.67\pm0.02} & \num{2.958\pm0.004} & \num{1.068}\\
    90.00 & 15 & 16:33:24 & \num{1.83\pm0.02} & \num{3.285\pm0.005} & \num{1.073}\\
    \bottomrule
  \end{tabular}
  \caption{[caption]%
    % Must say: per motor setting S, the number of photos, the time of the first (11 October 2026),
    % omega from the calibration with its uncertainty (common to all settings), the mean phase
    % shift of the photos relative to rest (ring minima; uncertainty: standard error of the
    % photos; the rest value's 0.02 rad is common to all rows and not included), and the ring
    % spacing P in r^2. Settings in the order taken.
  }
  \label{tab:settings}
\end{table}

% ---- Table 3, section 2.4 (optional: the methods; Fig. 7(b) and the text can carry them) ----
\begin{table}[tb]
  \centering\small\setlength{\tabcolsep}{4pt}
  \begin{tabular}{lccc}
    \toprule
    Method & $a$ & $s$ & $R^2$ \\
           & (\slopeunit) & (\unit{\radian}) & \\
    \midrule
    Ring minima        & \num{1.83\pm0.02} & 0.044 & 0.9987 \\
    Cosine fit         & \num{1.86\pm0.03} & 0.053 & 0.9982 \\
    Template matching  & \num{1.76\pm0.02} & 0.031 & 0.9993 \\
    \bottomrule
  \end{tabular}
  \caption{[caption]%
    % Must say: the slope from each method, with its fit uncertainty, the rms distance of the
    % settings from the line, and R^2. The method uncertainty is half the range of the three
    % slopes, 0.054 (2.9 %).
  }
  \label{tab:methods}
\end{table}

% ---- Table 4, section 2.5 (optional: the speed of light's budget) -------------------------
\begin{table}[tb]
  \centering\small
  \begin{tabular}{lc}
    \toprule
    Source & Relative uncertainty of $c$ \\
    \midrule
    Turn count $N$ & \qty{7.7}{\percent} \\
    Phase-reading method & \qty{2.9}{\percent} \\
    Loop radius $R$ (enters squared) & \qty{1.8}{\percent} \\
    Calibration & \qty{1.3}{\percent} \\
    Fit & \qty{1.3}{\percent} \\
    \midrule
    Total (in quadrature) & \qty{8.6}{\percent} \\
    \bottomrule
  \end{tabular}
  \caption{[caption]%
    % Must say: c = 8 pi^2 N R^2/(lambda a) = (2.9 +- 0.2) x 10^8 m/s; relative uncertainties added
    % in quadrature (c is proportional to N R^2/a); lambda given without an uncertainty (0.16 % per
    % nm, not included).
  }
  \label{tab:light}
\end{table}

% ---- Table 5, section 2.6 (uncertainty budget, full width) --------------------------------
\begin{table*}[tb]
  \centering\small
  \begin{tabular}{llll}
    \toprule
    Source & Kind & Size & Included \\
    \midrule
    Photo-to-photo variation & statistical & SEM \qtyrange{0.004}{0.028}{\radian} per setting & yes (in $s$) \\
    Wander of the phase between settings & statistical & about \qty{0.04}{\radian} per setting & yes (in $s$) \\
    Stopwatch: lap scatter and reaction time (\qty{0.2}{\second} per press) & statistical & \qtyrange{0.14}{0.9}{\percent} of a timed $\omega$ & yes (through the line) \\
    Calibration line and the table's scatter about it & statistical, correlated & \qty{1.3}{\percent} of $a$ & yes \\
    Phase-reading method & systematic & \qty{2.9}{\percent} of $a$ & yes \\
    Turn count $N = 13\pm1$ & given & \qty{7.7}{\percent} & yes ($a_\mathrm{th}$ and $c$) \\
    Loop radius $R = \qty{57.0\pm0.5}{\centi\metre}$ & given & \qty{1.8}{\percent} & yes ($a_\mathrm{th}$ and $c$) \\
    Wavelength $\lambda = \qty{640}{\nano\metre}$ & not given & \qty{0.16}{\percent} per nm & no \\
    Bend of the calibration at high settings & systematic & up to \qty{3}{\percent} of $\omega$ & partly (its rms is in the scatter) \\
    Steady drift in time & systematic & not separable from $a$ & no \\
    \bottomrule
  \end{tabular}
  \caption{[caption]%
    % Must say: every source considered, its kind, its size, and whether it is in the quoted
    % uncertainty. The slope's three parts (1.3 % fit, 2.9 % method, 1.3 % calibration) are added
    % in quadrature, 3.5 %: a = 1.83 +- 0.06. N and R enter only the prediction and c. Not checked:
    % the sign of the shift (rotation direction not recorded).
  }
  \label{tab:budget}
\end{table*}
